\section{What the sources establish and what is extended}
\label{sec:context}

\subsection{Pinned sources and the mathematical input}

The upstream source inspected here is \srcgrh, published on 25 November
2026. Its \code{grh/src/Statements.lean} imports Mathlib and states
nonvanishing for every Dirichlet $L$-function in $\Re s>1/2$, with the
principal pole excluded, as well as the corresponding critical-strip
zero statements. Its \code{Main.lean} also proves Mathlib's
\code{RiemannHypothesis}. These statements concern Dirichlet characters,
including imprimitive characters; they do not assert RH for arbitrary
automorphic or Artin $L$-functions \cite{Sernau2026}.

The repository pins Lean 4.35.0 and Mathlib commit
\code{56a005855ea17e0b0958f73ec801ae140b0a1679}.
Andrew Sutherland's independent verification report records successful
builds, axiom audits, statement comparisons and independent kernel checks
for the pinned upstream source \cite{Sutherland2026}.
The present work reads the source and those verification records; it does
not repeat the full upstream build. We use the stated Dirichlet-GRH theorem
as the analytic input of Section~\ref{grid:sec:main}.

The second source is the article at \srcarticle. Its README describes a
621-page consolidation of 103 reports and explicitly distinguishes that
consolidation from independent mathematical review \cite{ProveIt2026}.
We inspected its theorem and open-question structure, the relevant full
sections, and the three newest symplectic archives in \code{docs/incoming}.
This matters: the article's general symplectic-compression question has
already been treated in a newer incoming archive, so it is not presented
as a new contribution here.

\subsection{The source mechanism}

The upstream proof combines three inputs: a supply of short integral
forms, a spanning statement, and lower bounds forcing integral forms to
be long. In the determinant comparison, upper bounds for covolume and
index meet a lower bound obtained from twists and the product formula.
Cancellation of the common main quantity yields a one-sided estimate for
a weighted prime remainder. The meromorphic transform of that remainder
then excludes zeros to the right of the critical line. The precise entry
points are
\code{New/Reduction/MainInequality.lean},
\code{New/Reduction/OneSidedBound.lean}, and
\code{New/Reduction/Zeros/Criterion.lean} \cite{Sernau2026}.

The algebraic result below extends the global congruence term that enters
the denominator budget. The analytic result extends the subsequent
one-sided-transform and spectral-recovery interface. Thus the two results
use different parts of the supplied argument. Neither is obtained simply
by restating a standard prime-counting consequence of GRH.

\subsection{Precisely identified improvements}

\begin{longtable}{@{}p{0.24\textwidth}p{0.32\textwidth}p{0.36\textwidth}@{}}
\toprule
Interface & Inspected prior result & Result proved here\\
\midrule
Congruence modules &
\code{sel:thm:tracedepth}: length $\min(n,a)$ under full integral matrix
span; \code{sel:thm:bad}: extra torsion without it. &
The sharp bound $\min(n,a)+(d^2-1)\min(n,e)$ for nonmaximal image orders,
plus a bound on every invariant factor.\\[3pt]
Bad coefficient primes &
\code{sel:sec:open} requests an extension measured by the defect of the
image order. &
A finite conductor defect gives an explicit local and global
product-formula correction; extremal examples determine its coefficient.\\[3pt]
Finite prime sampling &
\code{ana:thm:binomial}: $(r+1)^2$ statistic samples, error
$O(\log r/\sqrt r)$, span $(1+\sqrt2)r$. &
$N$ ramp observations, error $O((\log N)^3/N^2)$, span $O(\log N)$,
with explicit dependence on fixed-family spectral constants.\\[3pt]
Central orders &
The subgroup construction recovers sums over characters, with individual
nonreal characters raised in \code{ana:q:family}. &
The real nonnegative weight $\Lambda(n)(\chi_0(n)+\Re\chi(n))$ recovers
the individual central order of $\chi$.\\
\bottomrule
\end{longtable}

The fourth entry is a simple useful corollary, not the main novelty claim.
The article also contains continuous convolution projectors and filters
using prescribed zero ordinates. The analytic advance is the finite-grid
rate, the control of all repeated aliases from a local count, and the
resulting sample and cutoff bounds. Uniform-average cascades are classical
digital filters \cite{Hogenauer1981}; no novelty is claimed for their
transfer polynomial. Likewise trace rigidity under residual absolute
irreducibility is classical \cite{Carayol1994,Brumer2019}; the new
algebraic assertion is the sharp quantitative bound without that residual
hypothesis.

Targeted searches of the literature and of the inspected sources did not
identify either exact main theorem. This is evidence for novelty relative
to the supplied program, not a proof of worldwide priority. The archive
records the search terms and the closest antecedents.
